% =====================================================================
% Relativistic Non-Dissipative Couette Flow with Magnetic Field
% Chapter IX, Sections 81--83: Relativistic Modification
% =====================================================================
% Requires: rel_preamble.tex (loaded by rel_main.tex)
%
% Classical reference: Chandrasekhar, Ch. IX, §§81--83
% Relativistic framework: Covariant ideal MHD, BDNK causality
% =====================================================================

\chapter*{Chapter IX (Relativistic): Non-Dissipative Couette Flow in
  Hydromagnetics}
\addcontentsline{toc}{chapter}{Chapter IX (Rel.): Non-Dissipative
  Couette Flow with Magnetic Field}

%======================================================================
\section{Relativistic MHD equations in cylindrical coordinates}
\label{sec:rel-9-80}
%======================================================================

In the classical theory (\S\,80), the equations governing an
incompressible, electrically conducting fluid in the presence of a
magnetic field are written in cylindrical coordinates $(r,\theta,z)$.
In the relativistic extension one must replace the Newtonian momentum
equation with the covariant conservation of the energy--momentum
tensor of a magnetised perfect fluid,
\begin{equation}
  \nabla_{\mu}\, T^{\mu\nu}_{\mathrm{MHD}} = 0,
  \label{eq:rel-9-Tcons}
\end{equation}
where (cf.\ the relativistic MHD framework, equation~\eqref{eq:Fmunu-decomp}
and RELATIVISTIC\_CONVENTIONS.md)
\begin{equation}
  T^{\mu\nu}_{\mathrm{MHD}}
  = \left(\varepsilon + p + \frac{b^{2}}{4\pi}\right)
    \frac{u^{\mu}\,u^{\nu}}{c^{2}}
  + \left(p + \frac{b^{2}}{8\pi}\right)g^{\mu\nu}
  - \frac{b^{\mu}\,b^{\nu}}{4\pi}.
  \label{eq:rel-9-TMHD}
\end{equation}
Here $\varepsilon$ is the total energy density (including rest mass),
$p$ the pressure, $u^{\mu}$ the 4-velocity normalised by
$u^{\mu}u_{\mu}=-c^{2}$, and $b^{\mu}$ the magnetic field 4-vector in
the fluid rest frame satisfying $b^{\mu}u_{\mu}=0$.  The relativistic
enthalpy density is
\begin{equation}
  w \;\equiv\; \frac{\varepsilon + p}{c^{2}}.
  \label{eq:rel-9-enthalpy}
\end{equation}

The ideal MHD condition (vanishing electric field in the fluid frame),
\begin{equation}
  F^{\mu\nu}\,u_{\nu} = 0,
  \label{eq:rel-9-idealMHD}
\end{equation}
together with the homogeneous Maxwell equations
$\nabla_{\mu}\,{}^{\ast}\!F^{\mu\nu}=0$, yields the relativistic
induction equation.  In Minkowski spacetime with cylindrical
coordinates $(t,r,\theta,z)$ and metric
$ds^{2} = -c^{2}\,dt^{2} + dr^{2} + r^{2}\,d\theta^{2} + dz^{2}$,
the spatial components of equation~\eqref{eq:rel-9-Tcons} and the
induction equation take the forms:

\medskip\noindent\textbf{Radial momentum:}
\begin{align}
  w\!\left(\frac{\partial u_{r}}{\partial t}
    + (\mathbf{u}\!\cdot\!\mathrm{grad})\,u_{r}
    - \frac{u_{\theta}^{2}}{r}\right)
  &- \frac{1}{4\pi}\!\left((\mathbf{b}\!\cdot\!\mathrm{grad})\,b_{r}
    - \frac{b_{\theta}^{2}}{r}\right) \notag\\
  &= -\frac{\partial}{\partial r}
      \!\left(p + \frac{b^{2}}{8\pi}\right)
    + \frac{w}{c^{2}}\,\mathcal{R}_{r},
  \label{eq:rel-9-1}
\end{align}
where $\mathcal{R}_{r}$ collects terms of order $u^{2}/c^{2}$
relative to the leading Newtonian terms.  Analogous equations hold for
the $\theta$ and $z$ components.

\medskip\noindent\textbf{Induction (frozen-in flux):}
\begin{equation}
  \frac{\partial b_{r}}{\partial t}
  + (\mathbf{u}\!\cdot\!\mathrm{grad})\,b_{r}
  - (\mathbf{b}\!\cdot\!\mathrm{grad})\,u_{r}
  = -\frac{b_{r}}{c^{2}}\,\frac{Du_{r}}{Dt}
  + O\!\left(\frac{u^{2}}{c^{2}}\right)\!b_{r},
  \label{eq:rel-9-induction-r}
\end{equation}
with $D/Dt$ the comoving (proper-time) derivative.  The key
relativistic modification is that the \emph{inertial mass density}
$\rho$ of the classical equations is everywhere replaced by the
\emph{relativistic enthalpy density}
\begin{equation}
  \rho \;\longrightarrow\; w
  = \frac{\varepsilon + p}{c^{2}}
  = \rho\!\left(1 + \frac{e + p/\rho}{c^{2}}\right),
  \label{eq:rel-9-rho-to-w}
\end{equation}
where $e$ is the specific internal energy.  The divergence-free
conditions
$\nabla_{\mu}(\rho\, u^{\mu})=0$ (baryon conservation) and
$\nabla_{\mu}\,b^{\mu}=0$ retain the classical cylindrical forms
\eqref{eq:9-12} and \eqref{eq:9-13} at leading order in $v/c$.

\medskip
The system admits the same stationary Couette solution as classically:
\begin{equation}
  u_{r} = u_{z} = 0,\quad u_{\theta} = V(r) = r\,\Omega(r),\quad
  b_{r} = b_{\theta} = 0,\quad b_{z} = B_{0} = \text{const},
  \label{eq:rel-9-baseflow}
\end{equation}
with the equilibrium condition
\begin{equation}
  \frac{d}{dr}\!\left(p + \frac{B_{0}^{2}}{8\pi}\right)
  = w\,\frac{V^{2}}{r}
  = \frac{(\varepsilon + p)}{c^{2}}\,\frac{V^{2}}{r}.
  \label{eq:rel-9-equil}
\end{equation}
Note that the centrifugal term involves $w = (\varepsilon+p)/c^{2}$
rather than $\rho$: this is the essential relativistic modification of
the radial force balance.

%======================================================================
\section{Stability of non-dissipative Couette flow with axial
  magnetic field: relativistic}
\label{sec:rel-81}
%======================================================================

We linearise about the base state~\eqref{eq:rel-9-baseflow} exactly as
in the classical \S\,81, with all perturbations proportional to
$\exp[i(\sigma t + m\theta + kz)]$ and $\sigma = p + m\Omega$.
The relativistic perturbation equations differ from the classical ones
(equations~\eqref{eq:9-24}--\eqref{eq:9-29}) by the replacement
$\rho \to w$ wherever $\rho$ appears as an inertial mass.  Defining
the \emph{relativistic Alfv\'en frequency}
\begin{equation}
  \Omega_{A,\mathrm{rel}}^{2}
  \;\equiv\; \frac{B_{0}^{2}\,k^{2}}{4\pi\,w}
  = \frac{B_{0}^{2}\,k^{2}\,c^{2}}{4\pi\,(\varepsilon + p)},
  \label{eq:rel-9-OmegaA}
\end{equation}
(compare the classical $\Omega_{A}^{2} = \mu H^{2}k^{2}/(4\pi\rho)$,
equation~\eqref{eq:9-43}), and introducing the Lagrangian displacement
$\boldsymbol{\xi}$ as in~\eqref{eq:9-38}, the perturbation equations
in terms of $\boldsymbol{\xi}$ take the identical form to
equations~\eqref{eq:9-40}--\eqref{eq:9-44}, but with:
\begin{enumerate}
\item $\Omega_{A}^{2}$ replaced by $\Omega_{A,\mathrm{rel}}^{2}$;
\item Rayleigh's discriminant generalised to
\begin{equation}
  \Phi_{\mathrm{rel}}(r)
  = \frac{1}{r^{3}}\,\frac{d}{dr}\!\left(r^{2}\,\Omega\right)^{2}
  = \frac{1}{r^{3}}\,\frac{d}{dr}\!\left(r^{2}\,V/r\right)^{2},
  \label{eq:rel-9-Phi}
\end{equation}
which is structurally unchanged, since the angular velocity profile
$\Omega(r)$ is unmodified;
\item the angular momentum per unit \emph{enthalpy}
\begin{equation}
  \ell_{\mathrm{rel}} = w\,r^{2}\,\Omega
  = \frac{(\varepsilon + p)}{c^{2}}\;r^{2}\,\Omega
  \label{eq:rel-9-angmom}
\end{equation}
replacing the classical $\rho\, r^{2}\Omega$ in the stability
criterion.
\end{enumerate}

The combined system~\eqref{eq:9-46} and \eqref{eq:9-49} becomes
\begin{equation}
  (\sigma^{2} - \Omega_{A,\mathrm{rel}}^{2})\,D_{*}\xi_{r}
  - \frac{2m\sigma\Omega}{r}\,\xi_{r}
  = \left(\frac{m^{2}}{r^{2}} + k^{2}\right)\varpi,
  \label{eq:rel-9-46}
\end{equation}
\begin{equation}
  \left(\sigma^{2} - \Omega_{A,\mathrm{rel}}^{2}
    - \Phi_{\mathrm{rel}}(r)
    - \frac{4\Omega^{2}\,\Omega_{A,\mathrm{rel}}^{2}}
           {\sigma^{2} - \Omega_{A,\mathrm{rel}}^{2}}\right)\xi_{r}
  = D\varpi
    + \frac{2m\sigma\Omega}{\sigma^{2} - \Omega_{A,\mathrm{rel}}^{2}}
      \,\frac{\varpi}{r},
  \label{eq:rel-9-49}
\end{equation}
with boundary conditions $\xi_{r}=0$ at $r=R_{1}$ and $R_{2}$.

%----------------------------------------------------------------------
\subsection*{(a) The case $m=0$: relativistic criterion}
\label{sec:rel-81a}
%----------------------------------------------------------------------

For axisymmetric perturbations ($m=0$), $\sigma = p$ and the system
reduces to the eigenvalue problem
\begin{equation}
  \kappa_{\mathrm{rel}}\,(DD_{*} - k^{2})\,\xi_{r}
  = -k^{2}\!\left[\Phi_{\mathrm{rel}}(r)
    + \frac{4\Omega^{2}\,\Omega_{A,\mathrm{rel}}^{2}}
           {\kappa_{\mathrm{rel}}}\right]\xi_{r},
  \label{eq:rel-9-53}
\end{equation}
where
\begin{equation}
  \kappa_{\mathrm{rel}} = p^{2} - \Omega_{A,\mathrm{rel}}^{2}.
  \label{eq:rel-9-54}
\end{equation}

The proof that the eigenvalues $\kappa_{\mathrm{rel}}$ are real
carries over \emph{verbatim} from the classical analysis
(\S\,81\,(a), equations~\eqref{eq:9-55}--\eqref{eq:9-57}), since the
self-adjointness of the operator depends only on the structure of the
equation, not on the numerical value of $\Omega_{A}$.

The integral identity~\eqref{eq:9-58} becomes
\begin{equation}
  \kappa_{\mathrm{rel}}^{2}\,I_{1}
  - \kappa_{\mathrm{rel}}\,k^{2}\,I_{2,\mathrm{rel}}
  - 4k^{2}\,\Omega_{A,\mathrm{rel}}^{2}\,I_{3} = 0,
  \label{eq:rel-9-58}
\end{equation}
with
\begin{equation}
  I_{2,\mathrm{rel}}
  = \int_{R_{1}}^{R_{2}} r\,\Phi_{\mathrm{rel}}(r)\,\xi_{r}^{2}\,dr.
  \label{eq:rel-9-I2}
\end{equation}
Since $\Phi_{\mathrm{rel}}(r)$ has the same functional form as
classically (only the Alfv\'en frequency differs), the variational
principle and the Sturmian structure of the spectrum are preserved.

\medskip\noindent\textbf{Stability criterion (relativistic).}
\emph{A necessary and sufficient condition for the stability of an
ideal MHD Couette flow with a uniform axial field, under axisymmetric
perturbations, is that}
\begin{equation}
  \boxed{%
  \frac{d}{dr}\!\left(\frac{(\varepsilon+p)}{c^{2}}\;
    r^{2}\,\Omega\right)^{2} > 0
  \quad\text{throughout } (R_{1},R_{2}).}
  \label{eq:rel-9-stability}
\end{equation}
This is the relativistic Rayleigh criterion: \emph{stability requires
that the specific angular momentum per unit enthalpy,}
$\ell_{\mathrm{rel}} = w\,r^{2}\Omega$, \emph{increase outward}.
In the non-relativistic limit $(\varepsilon+p)/c^{2}\to\rho$, the
classical Rayleigh criterion $d(\rho\,r^{2}\Omega)^{2}/dr > 0$ is
recovered.

\medskip
For a fluid with a relativistic equation of state (e.g.\ hot plasma
with $p\sim\varepsilon/3$), the enthalpy factor $w = (\varepsilon+p)/c^{2}$
can significantly exceed $\rho$, making the effective inertia larger
and the angular-momentum gradient steeper.  This relativistic
enhancement of inertia can either stabilise or destabilise the flow
depending on the radial profiles of $\varepsilon(r)$ and $p(r)$.

%----------------------------------------------------------------------
\subsection*{(b) Variational principle for $\lambda^{2}$: relativistic}
\label{sec:rel-81b}
%----------------------------------------------------------------------

The quadratic~\eqref{eq:rel-9-58} provides the basis for a variational
formulation exactly as in \S\,81\,(a)\,(ii).  Defining
\begin{equation}
  \lambda_{\mathrm{rel}}^{2}
  = \frac{\displaystyle\int_{R_{1}}^{R_{2}}
      (c_{\mathrm{rel}} - W)^{2}\,
      \bigl\{(D_{*}\chi)^{2} + k^{2}\chi^{2}\bigr\}\,r\,dr}
    {\displaystyle\int_{R_{1}}^{R_{2}}
      \phi_{\mathrm{rel}}(r)\,\chi^{2}\,r\,dr},
  \label{eq:rel-9-66}
\end{equation}
where $W = r\Omega$ and $\phi_{\mathrm{rel}}$ involves
$\Phi_{\mathrm{rel}}$ defined in~\eqref{eq:rel-9-Phi}, the stationary
values of $\lambda_{\mathrm{rel}}^{2}$ with respect to variations of
$\chi$ (subject to $\chi(R_{1})=\chi(R_{2})=0$) yield the eigenvalues.

The Sturmian monotonicity properties established classically
(equations~\eqref{eq:9-67}--\eqref{eq:9-74}) carry over to the
relativistic case without modification, since they rely on the
positivity structure of the integrals, which is unchanged.
Consequently, for each assigned $k^{2}$ and $\lambda_{\mathrm{rel}}^{2}>0$
there exist exactly two real characteristic values of the phase speed,
and no complex roots arise when $\Phi_{\mathrm{rel}}>0$ everywhere.

%======================================================================
\section{Oscillation periods with axial magnetic field: relativistic}
\label{sec:rel-82}
%======================================================================

%----------------------------------------------------------------------
\subsection*{(a) $\Omega = \mathrm{constant}$: relativistic
  magneto-inertial oscillations}
\label{sec:rel-82a}
%----------------------------------------------------------------------

When $\Omega = \text{const}$, the analysis of \S\,82\,(a) carries
through with $\Omega_{A}\to\Omega_{A,\mathrm{rel}}$.  The dispersion
relation~\eqref{eq:9-82} becomes
\begin{equation}
  \sigma^{2} - \sigma_{0}\,\sigma - \Omega_{A,\mathrm{rel}}^{2} = 0,
  \label{eq:rel-9-82}
\end{equation}
with roots
\begin{equation}
  \sigma = \tfrac{1}{2}\!\left[\sigma_{0}
    \pm \sqrt{\sigma_{0}^{2}
    + 4\,\Omega_{A,\mathrm{rel}}^{2}}\;\right].
  \label{eq:rel-9-84}
\end{equation}
Here $\sigma_{0}$ is the hydrodynamic eigenvalue from \S\,68\,(a).  The
splitting~\eqref{eq:9-83} generalises to
\begin{equation}
  \sigma_{1} + \sigma_{2} = \sigma_{0},
  \qquad
  \sigma_{1}\,\sigma_{2} = -\Omega_{A,\mathrm{rel}}^{2},
  \label{eq:rel-9-83}
\end{equation}
so that the magnetic field still splits each hydrodynamic mode into two
magneto-inertial oscillations, but the splitting magnitude now depends
on the relativistic Alfv\'en frequency
$\Omega_{A,\mathrm{rel}}^{2} = B_{0}^{2}k^{2}c^{2}/[4\pi(\varepsilon+p)]$
rather than $\mu H^{2}k^{2}/(4\pi\rho)$.

The oscillation period of the magneto-inertial modes is
\begin{equation}
  \mathcal{T}_{\mathrm{rel}}
  = \frac{2\pi}{|\sigma|}
  = \frac{2\pi}{\tfrac{1}{2}\bigl|\sigma_{0}
    \pm\sqrt{\sigma_{0}^{2}
    + 4\Omega_{A,\mathrm{rel}}^{2}}\;\bigr|}.
  \label{eq:rel-9-period}
\end{equation}
In the strong-field limit
$\Omega_{A,\mathrm{rel}}^{2} \gg \sigma_{0}^{2}$, the two roots
approach $\pm\Omega_{A,\mathrm{rel}}$ and the oscillation period
becomes that of a relativistic Alfv\'en wave:
\begin{equation}
  \mathcal{T}_{\mathrm{Alfven}}
  = \frac{2\pi}{\Omega_{A,\mathrm{rel}}}
  = \frac{2\pi}{k}\,
    \sqrt{\frac{4\pi\,(\varepsilon + p)}{B_{0}^{2}\,c^{2}}}.
  \label{eq:rel-9-TAlfven}
\end{equation}

%----------------------------------------------------------------------
\subsection*{(b) Stabilising effect of the magnetic field in narrow
  gaps: relativistic}
\label{sec:rel-82b}
%----------------------------------------------------------------------

For $m=0$ and $\Omega(r) = A + B/r^{2}$ (viscous Couette profile),
the narrow-gap analysis of \S\,82\,(b) carries through with the
replacement $\rho\to w = (\varepsilon+p)/c^{2}$.  The stability
criterion~\eqref{eq:9-96} generalises to
\begin{equation}
  \frac{B_{0}^{2}}{4\pi\,w}
  = \frac{B_{0}^{2}\,c^{2}}{4\pi\,(\varepsilon + p)}
  > 2\,\Omega_{1}^{2}\,R_{0}\,d\;
    \bigl[(1-\mu)\,\Lambda_{0}(\mu)\bigr],
  \label{eq:rel-9-96}
\end{equation}
where $\Lambda_{0}(\mu)$ is the same function tabulated in Table~XLI
of the classical treatment.  The criterion for complete stabilisation
of all adverse flows becomes
\begin{equation}
  \frac{B_{0}^{2}\,c^{2}}{4\pi\,(\varepsilon + p)}
  > q_{c}\;\bigl(2\,\Omega_{1}^{2}\,R_{0}\,d\bigr),
  \qquad q_{c} \simeq 0.0532.
  \label{eq:rel-9-100}
\end{equation}
Since $(\varepsilon+p)/c^{2} > \rho$, a \emph{stronger} magnetic field
is required in the relativistic regime to achieve the same
stabilising effect: the enhanced inertia demands a proportionally
larger magnetic tension.

%======================================================================
\section{Non-dissipative Couette flow with transverse (azimuthal)
  magnetic field: relativistic}
\label{sec:rel-83}
%======================================================================

When a current distribution parallel to the axis produces an
azimuthal magnetic field $B_{\theta}(r)$ (classical \S\,83), the
stationary solution~\eqref{eq:9-103} remains valid in the
relativistic theory with the equilibrium condition modified to
\begin{equation}
  \frac{d}{dr}\!\left(p + \frac{B_{\theta}^{2}}{8\pi}\right)
  = \frac{(\varepsilon + p)}{c^{2}}\,\frac{V^{2}}{r}
    - \frac{B_{\theta}^{2}}{4\pi\,r}.
  \label{eq:rel-9-104}
\end{equation}

The linearised equations for axisymmetric perturbations ($m=0$) lead,
as in the classical case, to the eigenvalue problem
\begin{equation}
  (DD_{*} - k^{2})\,\xi_{r}
  = -\frac{k^{2}}{p^{2}}\;
    \Psi_{\mathrm{rel}}(r)\,\xi_{r},
  \label{eq:rel-9-126}
\end{equation}
where the relativistic analogue of $\Psi(r)$
(equation~\eqref{eq:9-124}) is
\begin{equation}
  \Psi_{\mathrm{rel}}(r)
  = \frac{1}{w\,r^{3}}\,\frac{d}{dr}\bigl(w\,r^{2}\,\Omega\bigr)^{2}
    - \frac{r}{4\pi\,w}\,\frac{d}{dr}\!\left(\frac{B_{\theta}}{r}\right)^{\!2}.
  \label{eq:rel-9-Psi}
\end{equation}
The first term is the relativistic Rayleigh discriminant
$\Phi_{\mathrm{rel}}$ with enthalpy-weighted angular momentum; the
second accounts for the destabilising effect of the azimuthal magnetic
field gradient (the magnetic analogue of the centrifugal term).

\medskip\noindent\textbf{Michael's criterion (relativistic).}
\emph{A necessary and sufficient condition for stability of the flow
under axisymmetric perturbations is that
$\Psi_{\mathrm{rel}}(r) > 0$ throughout $(R_{1},R_{2})$.}  Explicitly,
\begin{equation}
  \boxed{%
  \frac{d}{dr}\!\left(\frac{\varepsilon+p}{c^{2}}\;
    r^{2}\,\Omega\right)^{\!2}
  > \frac{r^{4}}{4\pi}\,
    \frac{d}{dr}\!\left(\frac{B_{\theta}}{r}\right)^{\!2}
  \quad\text{throughout } (R_{1},R_{2}).}
  \label{eq:rel-9-Michael}
\end{equation}
In the non-relativistic limit $(\varepsilon+p)/c^{2}\to\rho$, this
recovers Michael's classical result.

For the non-axisymmetric modes ($m\ne 0$), the variational
formulation of \S\,83\,(b) extends straightforwardly with
$\Omega_{H}^{2}\to\Omega_{H,\mathrm{rel}}^{2}
= B_{\theta}^{2}/(4\pi\,w\,r^{2})$.  As in the classical theory, the
stability conditions for these modes have not been fully determined.

%======================================================================
\section{Causality constraints}
\label{sec:rel-9-causality}
%======================================================================

A crucial requirement of the relativistic theory, with no classical
counterpart, is that all characteristic speeds of the MHD system remain
subluminal.  We collect the relevant bounds.

\medskip\noindent\textbf{Alfv\'en speed.}
The relativistic Alfv\'en speed is
\begin{equation}
  v_{A}^{2}
  = \frac{b^{2}/4\pi}{w + b^{2}/(4\pi c^{2})}\;c^{2}
  = \frac{B_{0}^{2}\,c^{2}}{4\pi\,(\varepsilon + p) + B_{0}^{2}},
  \label{eq:rel-9-vA}
\end{equation}
which manifestly satisfies $v_{A} < c$ for any finite $B_{0}$.  (In
the classical limit $(\varepsilon+p)\to\rho c^{2} \gg B_{0}^{2}/4\pi$,
one recovers $v_{A}^{2}\to B_{0}^{2}/(4\pi\rho)$.)

\medskip\noindent\textbf{Magnetosonic speeds.}
The fast and slow magnetosonic speeds are (cf.\
RELATIVISTIC\_CONVENTIONS.md):
\begin{equation}
  v_{f,s}^{2}
  = \frac{1}{2}\!\left[c_{s}^{2} + v_{A}^{2}
    \pm\sqrt{(c_{s}^{2} + v_{A}^{2})^{2}
    - 4\,c_{s}^{2}\,v_{A}^{2}\cos^{2}\!\alpha}\;\right],
  \label{eq:rel-9-vfs}
\end{equation}
where $c_{s}^{2} = (\partial p/\partial\varepsilon)_{s}\le c^{2}$ is
the sound speed and $\alpha$ is the angle between the wave vector and
the magnetic field.  Both $v_{f}$ and $v_{s}$ satisfy $v_{f,s}\le c$,
as can be verified algebraically from $c_{s}\le c$ and $v_{A}\le c$.

\medskip\noindent\textbf{Implications for stability analysis.}
The eigenfrequencies $\sigma$ obtained from the
dispersion relation~\eqref{eq:rel-9-82} correspond to phase speeds
$v_{\mathrm{ph}} = \sigma/k$ that must satisfy
\begin{equation}
  |v_{\mathrm{ph}}| \le c.
  \label{eq:rel-9-subluminal}
\end{equation}
This is automatically guaranteed provided the Alfv\'en frequency is
defined through~\eqref{eq:rel-9-OmegaA} with the enthalpy density $w$
in the denominator, since
\begin{equation}
  \frac{\Omega_{A,\mathrm{rel}}^{2}}{k^{2}}
  = \frac{B_{0}^{2}\,c^{2}}{4\pi\,(\varepsilon + p)}
  < c^{2}
  \qquad\text{whenever } B_{0}^{2}/(4\pi) < \varepsilon + p,
  \label{eq:rel-9-causal-check}
\end{equation}
which holds for any physically reasonable equation of state (the
magnetic energy density is part of the total energy, hence always
bounded by $\varepsilon$).

For the magneto-inertial oscillations of \S\,\ref{sec:rel-82a}, the
group velocity can be computed from~\eqref{eq:rel-9-84}:
\begin{equation}
  v_{g} = \frac{\partial\sigma}{\partial k}
  = \frac{\Omega_{A,\mathrm{rel}}^{2}/k}
         {\sqrt{\sigma_{0}^{2}
           + 4\,\Omega_{A,\mathrm{rel}}^{2}}}.
  \label{eq:rel-9-vgroup}
\end{equation}
In the strong-field limit this approaches $v_{A} < c$, confirming
subluminal energy transport.

\medskip\noindent\textbf{Summary.}
The replacement $\rho\to w = (\varepsilon+p)/c^{2}$ in all inertial
terms is not merely a notational convenience but the mechanism by which
the relativistic theory enforces causality.  The denominator
$(\varepsilon+p)$ grows with the magnetic energy itself, preventing
the Alfv\'en speed and all MHD mode speeds from exceeding $c$.
This self-regulation is absent in Newtonian MHD, where arbitrarily
large $B$ fields produce arbitrarily large Alfv\'en speeds.

%======================================================================
\section{Astrophysical application: the MRI in Keplerian accretion disks}
\label{sec:rel-9-MRI-astro}
%======================================================================

The stability criteria of \S\S\,\ref{sec:rel-81}--\ref{sec:rel-83}
acquire profound astrophysical significance through the
\emph{magnetorotational instability} (MRI), identified by
Balbus \& Hawley~\cite{BalbusSHawley1991} as the driver of turbulence
in accretion disks.  The bounded Alfv\'en speed~\eqref{eq:rel-9-vA}
limits the unstable wavenumber band, while the maximum growth rate
$\sigma_{\max} = \frac{3}{4}|\Omega|$ is preserved for Keplerian
shear \cite{Balbus2003,PessahPsaltis2005}.

\begin{figure}[ht]
\centering
\includegraphics[width=0.95\textwidth]{fig_mri_keplerian.pdf}
\caption{MRI dispersion in a Keplerian disk: growth rate
  $\sigma/\Omega$ vs.\ wavenumber for several $v_A/c$.
  Left: classical.  Right: relativistic (bounded $v_A$).}
\label{fig:mri-keplerian}
\end{figure}

%======================================================================
\subsection{Application to Sgr~A$^{*}$: MRI wavelength vs.\ disk
  scale height}
\label{sec:rel-9-MRI-SgrA}
%======================================================================

As a concrete application of the relativistic MRI analysis, we
consider the accretion flow onto the Galactic Centre black hole,
Sgr~A$^{*}$, with $M = 4\times 10^{6}\,M_{\odot}$ and an
estimated accretion rate $\dot{M}\sim 10^{-8}\,M_{\odot}\,
\mathrm{yr}^{-1}$ \cite{Wielgus2020}.  At such low accretion
rates the flow is a radiatively inefficient accretion flow
(RIAF/ADAF; Narayan \& Yi~\cite{NarayanYi1994}), characterised
by ion temperatures approaching the virial value,
$T_{i}\sim 10^{12}\,(r_{g}/r)\;\mathrm{K}$, and geometrically
thick disks with $H/r\sim 0.3$--$1$.

The MRI characteristic wavelength is
\begin{equation}
  \lambda_{\mathrm{MRI}}
  = \frac{2\pi\,v_{A}}{\Omega_{K}}
  = \frac{2\pi}{k_{\max}},
  \label{eq:rel-9-lambdaMRI}
\end{equation}
where $v_{A}$ is the relativistic Alfv\'en speed~\eqref{eq:rel-9-vA}
and $k_{\max} = \Omega_{K}/v_{A}$ is the maximum unstable wavenumber.
A necessary condition for the MRI to operate within the disk is
\begin{equation}
  \lambda_{\mathrm{MRI}} \;\lesssim\; H(r),
  \label{eq:rel-9-resolved}
\end{equation}
i.e., the MRI must be \emph{resolved} by the disk vertical extent.
Equivalently, $v_{A}/(\Omega_{K}\,H) \lesssim 1/(2\pi)$.

For the ADAF model of Sgr~A$^{*}$, the disk density follows from
mass conservation,
$\rho \sim \dot{M}/(4\pi\,r\,H\,v_{r})$, with the radial inflow
velocity $v_{r}\sim\alpha\,c_{s}\,(H/r)^{2}$ and $\alpha\sim 0.1$.
The magnetic field strength is parametrised by the plasma beta,
$\beta_{\mathrm{pl}} = 8\pi p/B^{2}$.  Figure~\ref{fig:mri-sgra}
shows $\lambda_{\mathrm{MRI}}/H$ as a function of radius for three
values of $\beta_{\mathrm{pl}}$.

\begin{figure}[ht]
\centering
\includegraphics[width=0.95\textwidth]{fig_mri_sgra.pdf}
\caption{Left: ratio $\lambda_{\mathrm{MRI}}/H$ vs.\ $r/r_{g}$ for the
  Sgr~A$^{*}$ ADAF model with $M = 4\times 10^{6}\,M_{\odot}$,
  $\dot{M} = 10^{-8}\,M_{\odot}\,\mathrm{yr}^{-1}$.  The MRI is
  resolved ($\lambda_{\mathrm{MRI}} < H$) in the green shaded region.
  Right: disk aspect ratio $H/r$ and Alfv\'en speed $v_{A}/c$ vs.\
  radius.  Dashed lines show $v_{A}/c$ for three values of
  $\beta_{\mathrm{pl}}$.  In the strongly magnetised case
  ($\beta_{\mathrm{pl}} = 1$), $v_{A}/c$ can reach $\sim 10^{-2}$
  near the ISCO, making the relativistic Alfv\'en
  bound~\eqref{eq:rel-9-vA} relevant.}
\label{fig:mri-sgra}
\end{figure}

Key findings:
\begin{enumerate}
\item For $\beta_{\mathrm{pl}} \gtrsim 10$, the MRI is well resolved
  ($\lambda_{\mathrm{MRI}}/H \ll 1$) at all radii beyond the ISCO,
  confirming that turbulent angular momentum transport operates
  throughout the ADAF.
\item For $\beta_{\mathrm{pl}} \sim 1$ (magnetically dominated), the
  MRI wavelength can exceed $H$ in the inner regions ($r \lesssim
  10\,r_{g}$), where the field is dynamically important and the MRI
  enters the ``magnetically arrested'' regime
  (Liska et al.~\cite{Liska2021}).
\item The relativistic correction to $v_{A}$
  (equation~\eqref{eq:rel-9-vA}) is at most a few percent for
  Sgr~A$^{*}$, but becomes significant for higher-accretion-rate
  systems where $v_{A}/c \gtrsim 0.1$.
\end{enumerate}
These results connect the linear stability analysis of
\S\S\,\ref{sec:rel-81}--\ref{sec:rel-83} to the EHT observations of
M87$^{*}$ and Sgr~A$^{*}$ morphology (Wielgus et
al.~\cite{Wielgus2020}), where the turbulent accretion flow structure
is shaped by MRI-driven angular momentum transport operating under
the relativistic constraints derived here.

\medskip
\noindent\textbf{References for \S\S\,81--83.}
Classical: Chandrasekhar~\cite{Chandrasekhar1961}.
MRI: Balbus \& Hawley~\cite{BalbusSHawley1991},
Balbus~\cite{Balbus2003}.
Relativistic MRI: Pessah \& Psaltis~\cite{PessahPsaltis2005},
Pessah, Chan \& Psaltis~\cite{PessahChanPsaltis2008}.
Relativistic MHD: Rezzolla \& Zanotti~\cite{RezzollaZanotti2013}.
ADAF model: Narayan \& Yi~\cite{NarayanYi1994}.
EHT morphology: Wielgus et al.~\cite{Wielgus2020}.
GRMHD dynamo: Liska et al.~\cite{Liska2021}.

%======================================================================
\section*{Non-relativistic limit}
\label{sec:rel-9-NRlimit}
%======================================================================

Throughout this chapter the speed of light $c$ has been kept explicit.
The non-relativistic limit is obtained by taking
$\varepsilon\to\rho c^{2}$, $p\ll\rho c^{2}$, so that
$w\to\rho$, $\Omega_{A,\mathrm{rel}}\to\Omega_{A}$, and all
equations reduce to the classical forms of \S\S\,81--83.  The
causality bounds become vacuous ($v_{A}/c\to 0$), and the Rayleigh and
Michael criteria return to their standard Newtonian forms.

%% ================================================================
\section{Stability of multicomponent Israel--Stewart--Maxwell theory}
\label{sec:rel-9-multicomponent-IS}
%% ================================================================

The MHD stability analysis of the preceding sections treats the fluid
as a single-component system.  In realistic astrophysical plasmas (neutron
star cores, QGP at finite baryon density, heavy-ion collisions at lower
energies), multiple conserved charges coexist (baryon number, electric
charge, strangeness), each with its own diffusive current.
Gavassino \& Shokri (\citealt{GavassinoShokri2024}, arXiv:2307.11615)
derive the stability conditions for multicomponent Israel--Stewart
hydrodynamics in the presence of electromagnetic fields.

\subsection{The multicomponent IS-Maxwell model}

The fluid state is characterised by $\{\varepsilon, u^\mu, n_A,
V_A^\mu\}$, where $A = 1,\ldots,l$ labels the chemical species.
The entropy current includes a quadratic cost for diffusive currents:
\begin{equation}
\label{eq:rel-9-IS-entropy}
S^\mu = s\,u^\mu - \frac{\mu^A}{T}\,V_A^\mu
  - \frac{b^{AB}}{2T}\,V_A^\nu V_{B\nu}\,u^\mu\,,
\end{equation}
where $b^{AB}$ is a symmetric positive-definite matrix quantifying the
out-of-equilibrium entropy cost.  The key result is derived using the
information current method:

\begin{equation}
\label{eq:rel-9-IS-info-current}
T E^\mu = \biggl[\delta T\,\delta s + \delta\mu^A\,\delta n_A
  + (\varepsilon{+}P)\,\delta u^\nu\delta u_\nu
  + b^{AB}\,\delta V_A^\nu\delta V_{B\nu}\biggr]
  \frac{u^\mu}{2} + \cdots
\end{equation}
Stability requires $E^\mu$ to be future-directed non-spacelike for all
perturbations, which yields:
\begin{enumerate}
\item \textbf{Fluid sector}: the standard thermodynamic stability
  conditions ($\partial\varepsilon/\partial T > 0$, positive
  compressibility) must hold.
\item \textbf{Diffusive sector}: the matrix $b^{AB}$ must be positive
  definite.
\item \textbf{Coupling}: these conditions are \emph{independent}---the
  fluid and diffusive sectors decouple in the stability analysis.
\end{enumerate}

\subsection{Extension to dynamic electromagnetic fields}

The central result of Gavassino \& Shokri is that the electromagnetic
part of the information current is \emph{automatically} stable and causal:
\begin{quote}
\textit{The stability conditions of a dissipative hydrodynamic theory
extend unchanged to the corresponding theory of resistive dissipative
magnetohydrodynamics, provided the medium is non-polarizable and
non-magnetizable.}
\end{quote}
This means that the MRI stability analysis of
\S\S\,\ref{sec:rel-81}--\ref{sec:rel-83}, which uses ideal MHD, remains
valid when extended to resistive multicomponent plasmas: the diffusion
currents and electromagnetic fields do not introduce new instabilities
beyond those already present in the single-component inviscid theory,
as long as the thermodynamic stability conditions and the positivity of
$b^{AB}$ are satisfied.

\subsection{Connection to Carter's multifluid theory}

Gavassino \& Shokri show that the information current
\eqref{eq:rel-9-IS-info-current} can be mapped, via a change of
variables, to the information current of Carter's multifluid theory.
This establishes a formal equivalence between the IS and Carter approaches
for the purpose of stability analysis.  Gavassino (\citealt{GavassinoCarter2022},
arXiv:2202.06760) had previously derived the stability conditions for
Carter's multifluid theory itself, proving that:
\begin{enumerate}
\item The entrainment matrix $\mathcal{K}^{XY}$ (coupling the momenta
  of different species to the currents of other species) must be
  positive definite for stability.
\item Thermodynamic stability implies linear causality: if the entropy
  is maximised at equilibrium, all characteristic speeds are subluminal.
\end{enumerate}
These results are directly relevant to the superfluid Couette flows
of Chapter~VI, where the entrainment between neutron and proton
superfluids determines the stability of differential rotation in
neutron star interiors.
